\chapter{T\'itulo do primeiro cap\'itulo} \label{cap:cap1}

Escrever bem é uma arte que exige muita técnica e dedicação. Há vários bons livros sobre como escrever uma boa dissertação ou tese. Para a escrita de textos em Ciência da Computação, Writing for Computer Science \cite{zobel2014} é uma leitura obrigatória. O livro Metodologia de Pesquisa para Ciência da Computação \cite{wazlawick2009} também merece uma boa lida. 

\section{Introdu\c{c}\~ao}\label{sec:intro}

Apresentar uma visão geral do assunto que será abordado no trabalho, procurando fazer com que o leitor adquira uma compreensão inicial do que será tratado e fornecendo informações que o levem a perceber a sua importância.  

Exemplo de Figura: Ver Figura~\ref{fig:exefig}.

\begin{figure}[!ht]
\centering
\includegraphics[width=0.3\linewidth]{figuras/exefig.eps}
\caption{Exemplo de figura}
\label{fig:exefig}
\end{figure}

Este \'e o primeiro cap\'itulo. 

Exemplo de refer\^encias~\cite{tese2011, confinter2011, confnac2011, rev2011, site2011}.

Este \'e o primeiro cap\'itulo. 

Exemplo de Tabela: ver Tabela~\ref{tab:exetab}.

\begin{table}[!ht]
\begin{center}
\caption{Distribuição IMC em Adultos}
\label{tab:exetab}
\begin{tabular}{|c |c |}
\hline
\textbf{\textbf{Classificação}} & \textbf{IMC}\\
\hline\hline
Baixo Peso & < 18,5 \\
Peso Adequado & > 18,5 e < 25 \\
Sobrepeso & > 25 e < 30 \\
Obesidade & > 30 \\
\hline
\end{tabular}
\end{center}
\end{table}

Este \'e o primeiro cap\'itulo. 
 


\section{Desenvolvimento}\label{sec:desenv}

Conforme explicado na Se\c{c}\~ao~\ref{sec:intro}, nesta segunda se\c{c}\~ao ser\'a descrito o desenvolvimento.
